\subsection{DIFFERENTIAL EQUATIONS ADMITTING SOLUTIONS THAT ARE PRIMITIVES OF REGULATED FUNCTIONS}

Recall (II, p.~54, def. 3) that a vector function \(\mathbf{u}\) defined on an interval \(\mathrm{I} \subset \mathbf{R}\) is said to be regulated if it is the uniform limit of step functions on every compact subset of I; an equivalent condition is that at every interior point of I the function \(\mathbf{u}\) has a right and a left limit, and also a right limit at the left-hand end point of I and a left limit at the right-hand endpoint of I, when these two points belong to I (II, p.~54, th. 3). In this chapter we shall restrict ourselves to differential equations (1) for which every solution is a primitive of a regulated function on I. This condition is clearly satisfied if, for every continuous map \(\mathbf{u}\) of I into H, the function \(\mathbf{f}(t, \mathbf{u}(t))\) is regulated on I; the following lemma gives a sufficient condition for this:

Lemma 1. Let \(\mathbf{f}\) be a map from \(\mathrm{I} \times \mathrm{H}\) into \(\mathrm{E}\) such that, on writing \(\mathbf{f}_{\mathbf{x}}\) (for every \(\mathbf{x} \in \mathrm{H}\) ) for the map \(t \mapsto \mathbf{f}(t, \mathbf{x})\) of \(\mathrm{I}\) into \(\mathrm{E}\) , the following conditions are satisfied: \(1^{\circ}\) \(\mathbf{f}_{\mathbf{x}}\) is regulated on \(\mathrm{I}\) for every \(\mathbf{x} \in \mathrm{H}\) ; \(2^{\circ}\) the map \(\mathbf{x} \mapsto \mathbf{f}_{\mathbf{x}}\) of \(\mathrm{H}\) into the set \(\mathcal{F}(\mathrm{I}, \mathrm{E})\) of maps from \(\mathrm{I}\) into \(\mathrm{E}\) is continuous when one endows \(\mathcal{F}(\mathrm{I}, \mathrm{E})\) with the topology of compact convergence (Gen.~Top., X, p.~278). Under these conditions:

\(1^{\circ}\) For every continuous map \(\mathbf{u}\) of I into H the function \(t \mapsto \mathbf{f}(t, \mathbf{u}(t))\) is regulated on I; more precisely, the right (resp. left) limit of this function at a point \(t_0 \in \mathbf{I}\) is equal to the right (resp. left) limit of the function \(t \mapsto \mathbf{f}(t, \mathbf{u}(t_0))\) at the point \(t_0\) .

\(2^{\circ}\) If \((\mathbf{u}_n)\) is a sequence of maps of I into H which converges uniformly to a continuous function \(\mathbf{u}\) of I into H on every compact subset of I, then the sequence of functions \(t \mapsto \mathbf{f}(t, \mathbf{u}_n(t))\) converges uniformly to \(\mathbf{f}(t, \mathbf{u}(t))\) on every compact subset of I.

\(1^{\circ}\) Let \(\mathbf{c}\) be the right limit of \(\mathbf{f}(t, \mathbf{u}(t_0))\) at the point \(t_0\) ; for every \(\varepsilon > 0\) there is a compact neighbourhood V of \(t_0\) in I such that \(\| \mathbf{f}(t, \mathbf{u}(t_0)) - \mathbf{c} \| \leqslant \varepsilon\) for \(t \in V\) and \(t > t_0\) . On the other hand, there exists \(\delta > 0\) such that the relations

\[
\mathbf {x} \in \mathrm{H}, \quad \| \mathbf {x} - \mathbf {u} (t _ {0}) \| \leqslant \delta
\]

imply \(\|\mathbf{f}(s,\mathbf{x})-\mathbf{f}(s,\mathbf{u}(t_{0}))\|\leqslant\varepsilon\) for all \(s\in V\) ; if \(W\subset V\) is a neighbourhood of \(t_{0}\) in I such that \(\|\mathbf{u}(t)-\mathbf{u}(t_{0})\|\leqslant\delta\) for every \(t\in W\) , then \(\|\mathbf{f}(t,\mathbf{u}(t))-\mathbf{c}\|\leqslant2\varepsilon\) for \(t\in W\) and \(t>t_{0}\) , which proves that c is the right limit of \(\mathbf{f}(t,\mathbf{u}(t))\) at the point \(t_{0}\) .

\(2^{\circ}\) Let \(\mathbf{K}\) be a compact subset of \(\mathrm{I}\) ; since \(\mathbf{u}\) is continuous on \(\mathrm{I}, \mathbf{u}(\mathbf{K})\) is a compact subset of \(\mathrm{H}\) ; for every \(\varepsilon > 0\) and \(\mathbf{x} \in \mathbf{u}(\mathbf{K})\) there exists a number \(\delta_{\mathbf{x}}\) such that, for every \(\mathbf{y} \in \mathrm{H}\) , \(\| \mathbf{y} - \mathbf{x} \| \leqslant \delta_{\mathbf{x}}\) and every \(t \in \mathbf{K}\) , one has \(\| \mathbf{f}(t, \mathbf{y}) - \mathbf{f}(t, \mathbf{x}) \| \leqslant \varepsilon\) . There is a finite number of points \(\mathbf{x}_i \in \mathbf{u}(\mathbf{K})\) such that the closed balls with centre \(\mathbf{x}_i\) and radius \(\frac{1}{2} \delta_{\mathbf{x}_i}\) form a cover of \(\mathbf{u}(\mathbf{K})\) ; let \(\delta = \operatorname{Min}\left(\delta_{\mathbf{x}_i}\right)\) . By hypothesis there exists an integer \(n_0\) such that for \(n \geqslant n_0\) one has \(\| \mathbf{u}_n(t) - \mathbf{u}(t) \| \leqslant \frac{1}{2} \delta\) for every \(t \in \mathbf{K}\) . Now, for every \(t \in \mathbf{K}\) there exists an index \(i\) such that

\[
\left\| \mathbf {u} (t) - \mathbf {x} _ {i} \right\| \leqslant \frac {1}{2} \delta_ {\mathbf {x} _ {i}};
\]

consequently one has \(\|\mathbf{u}_{n}(t)-\mathbf{x}_{i}\|\leqslant\delta_{\mathbf{x}}\) , whence \(\|\mathbf{f}(t,\mathbf{u}_{n}(t))-\mathbf{f}(t,\mathbf{u}(t))\|\leqslant2\varepsilon\) for every \(t\in K\) and every \(n\geqslant n_{0}\) .

For the rest of this section I will denote an interval contained in R, not reducing to a single point, H an open nonempty set contained in the normed space E, and f a map from \(I \times H\) into E satisfying the hypotheses of lemma 1.

PROPOSITION 1. Let \(t_0\) be a point of I and \(\mathbf{x}_0\) a point of H; for a continuous function \(\mathbf{u}\) to be a solution of the equation (1) on I and take the value \(\mathbf{x}_0\) at the point \(t_0\) , it is necessary and sufficient that it satisfies the relation

\[
\mathbf {u} (t) = \mathbf {x} _ {0} + \int_ {t _ {0}} ^ {t} \mathbf {f} (s, \mathbf {u} (s)) d s\tag{6}
\]

for every \(t\in \mathrm{I}\)

Indeed, by lemma 1, if \(\mathbf{u}\) is a solution of (1) on I, then \(\mathbf{f}(t, \mathbf{u}(t))\) is regulated, so the right-hand side of (6) is defined and equal to \(\mathbf{u}(t)\) for every \(t \in I\) . Conversely, if \(\mathbf{u}\) is a continuous function that satisfies (6) then \(\mathbf{f}(t, \mathbf{u}(t))\) is regulated, by lemma 1, so \(\mathbf{u}\) has derivative equal to \(\mathbf{f}(t, \mathbf{u}(t))\) except at the points of a countable subset of I.

COROLLARY. At every point of I distinct from the left (resp. right) endpoint of this interval, every solution u of (1) on I admits a left (resp. right) derivative equal to the left (right) limit of \(\mathbf{f}(t,\mathbf{u}(t))\) at this point.

PROPOSITION 2. If \(\mathbf{f}\) is a continuous map from \(\mathrm{I} \times \mathrm{H}\) into \(\mathrm{E}\) , then every solution of (1) on \(\mathrm{I}\) is a strict solution.

Indeed, such a solution \(\mathbf{u}\) is a primitive of the continuous function \(\mathbf{f}(t,\mathbf{u}(t))\) (II, p.~66, prop. 3).

Furthermore, we note that a continuous function \(\mathbf{f}\) on \(\mathrm{I} \times \mathrm{H}\) satisfies the conditions of lemma 1 (Gen.~Top., X, p.~286, cor. 3).

In the sequel we shall choose \(t_{0} \in I\) and \(x_{0} \in H\) arbitrarily and investigate whether there exist solutions of (1) on I (or on a neighbourhood of \(t_{0}\) in I) taking the value \(x_{0}\) at the point \(t_{0}\) (or, what comes to the same, solutions of (6)).

